% Secciones 3 a 6: GLC, análisis, FIRST/FOLLOW, tabla LL(1) (tarea R1) - autor: Persona B.
\section{Gramática libre de contexto}

La gramática de DomoDSL es $G = (N, T, P, S)$ con:
\begin{itemize}[nosep]
\item \textbf{No terminales} $N$ (14): \textit{programa}, \textit{programa}$'$, \textit{regla}, \textit{condicion}, \textit{condicion}$'$, \textit{termino}, \textit{termino}$'$, \textit{factor}, \textit{comparacion}, \textit{valor}, \textit{acciones}, \textit{acciones}$'$, \textit{accion}, \textit{parametros}.
\item \textbf{Terminales} $T$ (14, son los tokens): \texttt{WHEN}, \texttt{THEN}, \texttt{PUNTO\_COMA}, \texttt{OR}, \texttt{AND}, \texttt{NOT}, \texttt{PARENT\_IZQ}, \texttt{PARENT\_DER}, \texttt{ID}, \texttt{OP\_REL}, \texttt{NUM}, \texttt{TRUE}, \texttt{FALSE}, \texttt{COMA}.
\item \textbf{Símbolo inicial} $S$ = \textit{programa}.
\item \textbf{Producciones} $P$: las 24 que muestra la tabla siguiente. (Las producciones numeradas son las que usa el
código; el símbolo \textit{x}$'$ se escribe \texttt{x\_p} en el programa.)
\end{itemize}

\begin{longtable}{@{}r p{11cm}@{}}
\toprule
\textbf{N.º} & \textbf{Producción} \\
\midrule
\endhead
1 & \textit{programa} $\rightarrow$ \textit{regla} \textit{programa}$'$ \\
2 & \textit{programa}$'$ $\rightarrow$ \textit{regla} \textit{programa}$'$ \\
3 & \textit{programa}$'$ $\rightarrow$ $\varepsilon$ \\
4 & \textit{regla} $\rightarrow$ \texttt{WHEN} \textit{condicion} \texttt{THEN} \textit{acciones} \texttt{PUNTO\_COMA} \\
5 & \textit{condicion} $\rightarrow$ \textit{termino} \textit{condicion}$'$ \\
6 & \textit{condicion}$'$ $\rightarrow$ \texttt{OR} \textit{termino} \textit{condicion}$'$ \\
7 & \textit{condicion}$'$ $\rightarrow$ $\varepsilon$ \\
8 & \textit{termino} $\rightarrow$ \textit{factor} \textit{termino}$'$ \\
9 & \textit{termino}$'$ $\rightarrow$ \texttt{AND} \textit{factor} \textit{termino}$'$ \\
10 & \textit{termino}$'$ $\rightarrow$ $\varepsilon$ \\
11 & \textit{factor} $\rightarrow$ \texttt{NOT} \textit{factor} \\
12 & \textit{factor} $\rightarrow$ \texttt{PARENT\_IZQ} \textit{condicion} \texttt{PARENT\_DER} \\
13 & \textit{factor} $\rightarrow$ \textit{comparacion} \\
14 & \textit{comparacion} $\rightarrow$ \texttt{ID} \texttt{OP\_REL} \textit{valor} \\
15 & \textit{valor} $\rightarrow$ \texttt{NUM} \\
16 & \textit{valor} $\rightarrow$ \texttt{TRUE} \\
17 & \textit{valor} $\rightarrow$ \texttt{FALSE} \\
18 & \textit{valor} $\rightarrow$ \texttt{ID} \\
19 & \textit{acciones} $\rightarrow$ \textit{accion} \textit{acciones}$'$ \\
20 & \textit{acciones}$'$ $\rightarrow$ \texttt{COMA} \textit{accion} \textit{acciones}$'$ \\
21 & \textit{acciones}$'$ $\rightarrow$ $\varepsilon$ \\
22 & \textit{accion} $\rightarrow$ \texttt{ID} \texttt{ID} \textit{parametros} \\
23 & \textit{parametros} $\rightarrow$ \textit{valor} \textit{parametros} \\
24 & \textit{parametros} $\rightarrow$ $\varepsilon$ \\
\bottomrule
\end{longtable}

\paragraph{Lectura de la gramática.}
\begin{itemize}[nosep]
\item Un \textbf{programa} es una o más reglas (1--3). Una \textbf{regla} es \texttt{WHEN condición THEN acciones ;} (4).
\item Las \textbf{condiciones} se estratifican en tres niveles para codificar la precedencia: \textit{condicion} (OR, la menor),
\textit{termino} (AND) y \textit{factor} (NOT, paréntesis y comparaciones, la mayor) (5--13). Una \textbf{comparación} es
\texttt{ID OP\_REL valor} (14) y un \textbf{valor} es un número, un booleano o un identificador (15--18).
\item Las \textbf{acciones} son una lista separada por comas (19--21); cada una es \texttt{verbo dispositivo [parámetros]} (22--24), donde los
parámetros son cero o más valores.
\end{itemize}

\paragraph{Ejemplo de derivación.} Derivación por la izquierda de \texttt{WHEN a > 1 THEN x y;} (los terminales se muestran como
tipos de token: \texttt{ID}, \texttt{OP\_REL}, \texttt{NUM}, \dots). La derivación fue reconstruida a partir de la traza del analizador LL(1).

\begin{longtable}{@{}r r >{\raggedright\arraybackslash}p{12.2cm}@{}}
\toprule
\textbf{Paso} & \textbf{Prod.} & \textbf{Forma sentencial (derivación por la izquierda)} \\
\midrule
\endhead
0 &  & \footnotesize \textit{programa} \\
1 & 1 & $\Rightarrow$ \footnotesize \textit{regla} \textit{programa}$'$ \\
2 & 4 & $\Rightarrow$ \footnotesize \texttt{WHEN} \textit{condicion} \texttt{THEN} \textit{acciones} \texttt{PUNTO\_COMA} \textit{programa}$'$ \\
3 & 5 & $\Rightarrow$ \footnotesize \texttt{WHEN} \textit{termino} \textit{condicion}$'$ \texttt{THEN} \textit{acciones} \texttt{PUNTO\_COMA} \textit{programa}$'$ \\
4 & 8 & $\Rightarrow$ \footnotesize \texttt{WHEN} \textit{factor} \textit{termino}$'$ \textit{condicion}$'$ \texttt{THEN} \textit{acciones} \texttt{PUNTO\_COMA} \textit{programa}$'$ \\
5 & 13 & $\Rightarrow$ \footnotesize \texttt{WHEN} \textit{comparacion} \textit{termino}$'$ \textit{condicion}$'$ \texttt{THEN} \textit{acciones} \texttt{PUNTO\_COMA} \textit{programa}$'$ \\
6 & 14 & $\Rightarrow$ \footnotesize \texttt{WHEN} \texttt{ID} \texttt{OP\_REL} \textit{valor} \textit{termino}$'$ \textit{condicion}$'$ \texttt{THEN} \textit{acciones} \texttt{PUNTO\_COMA} \textit{programa}$'$ \\
7 & 15 & $\Rightarrow$ \footnotesize \texttt{WHEN} \texttt{ID} \texttt{OP\_REL} \texttt{NUM} \textit{termino}$'$ \textit{condicion}$'$ \texttt{THEN} \textit{acciones} \texttt{PUNTO\_COMA} \textit{programa}$'$ \\
8 & 10 & $\Rightarrow$ \footnotesize \texttt{WHEN} \texttt{ID} \texttt{OP\_REL} \texttt{NUM} \textit{condicion}$'$ \texttt{THEN} \textit{acciones} \texttt{PUNTO\_COMA} \textit{programa}$'$ \\
9 & 7 & $\Rightarrow$ \footnotesize \texttt{WHEN} \texttt{ID} \texttt{OP\_REL} \texttt{NUM} \texttt{THEN} \textit{acciones} \texttt{PUNTO\_COMA} \textit{programa}$'$ \\
10 & 19 & $\Rightarrow$ \footnotesize \texttt{WHEN} \texttt{ID} \texttt{OP\_REL} \texttt{NUM} \texttt{THEN} \textit{accion} \textit{acciones}$'$ \texttt{PUNTO\_COMA} \textit{programa}$'$ \\
11 & 22 & $\Rightarrow$ \footnotesize \texttt{WHEN} \texttt{ID} \texttt{OP\_REL} \texttt{NUM} \texttt{THEN} \texttt{ID} \texttt{ID} \textit{parametros} \textit{acciones}$'$ \texttt{PUNTO\_COMA} \textit{programa}$'$ \\
12 & 24 & $\Rightarrow$ \footnotesize \texttt{WHEN} \texttt{ID} \texttt{OP\_REL} \texttt{NUM} \texttt{THEN} \texttt{ID} \texttt{ID} \textit{acciones}$'$ \texttt{PUNTO\_COMA} \textit{programa}$'$ \\
13 & 21 & $\Rightarrow$ \footnotesize \texttt{WHEN} \texttt{ID} \texttt{OP\_REL} \texttt{NUM} \texttt{THEN} \texttt{ID} \texttt{ID} \texttt{PUNTO\_COMA} \textit{programa}$'$ \\
14 & 3 & $\Rightarrow$ \footnotesize \texttt{WHEN} \texttt{ID} \texttt{OP\_REL} \texttt{NUM} \texttt{THEN} \texttt{ID} \texttt{ID} \texttt{PUNTO\_COMA} \\
\bottomrule
\end{longtable}

\section{Análisis de la gramática}

\subsection{Terminales, no terminales y símbolo inicial}
Ya identificados en la sección~3: 14 terminales (los tokens), 14 no terminales y símbolo inicial \textit{programa}.
Los no terminales \textit{programa'}, \textit{condicion'}, \textit{termino'}, \textit{acciones'} son \emph{auxiliares}
introducidos al transformar la gramática; el resto tiene significado propio del dominio.

\subsection{Primera versión de la gramática ($G_0$)}
Al escribir la gramática de forma natural se obtiene:
\begin{lstlisting}[language={},numbers=none]
programa    → regla programa | regla
regla       → WHEN condicion THEN acciones PUNTO_COMA
condicion   → condicion OR condicion | condicion AND condicion
            | NOT condicion | PARENT_IZQ condicion PARENT_DER | comparacion
comparacion → ID OP_REL valor
valor       → NUM | TRUE | FALSE | ID
acciones    → acciones COMA accion | accion
accion      → ID ID parametros
parametros  → parametros valor | ε
\end{lstlisting}
$G_0$ describe el lenguaje deseado, pero \textbf{no sirve para un analizador predictivo}: tiene recursión por la
izquierda, prefijos comunes y es ambigua. A continuación se analiza cada problema.

\subsection{Recursión por la izquierda}
En $G_0$ hay recursión izquierda \emph{directa} en tres no terminales: \textit{condicion} ($\textit{condicion}\to\textit{condicion}\ \texttt{OR}\ \dots$),
\textit{acciones} y \textit{parametros}. Un parser de descenso recursivo entraría en bucle infinito. Se aplica la
transformación estándar $A\to A\alpha\mid\beta \;\Longrightarrow\; A\to\beta A',\ A'\to\alpha A'\mid\varepsilon$:

\begin{center}\small
\begin{tabular}{@{}>{\raggedright\arraybackslash}p{5.2cm} >{\raggedright\arraybackslash}p{7.3cm} >{\raggedright\arraybackslash}p{2.6cm}@{}}
\toprule
\textbf{Original ($G_0$)} & \textbf{Transformado} & \textbf{En la gramática final} \\
\midrule
$\textit{acciones}\to\textit{acciones}\ \texttt{COMA}\ \textit{accion}\mid\textit{accion}$ &
$\textit{acciones}\to\textit{accion}\ \textit{acciones'}$ & prods. 19--21 \\
 & $\textit{acciones'}\to\texttt{COMA}\ \textit{accion}\ \textit{acciones'}\mid\varepsilon$ & \\
$\textit{parametros}\to\textit{parametros}\ \textit{valor}\mid\varepsilon$ &
$\textit{parametros}\to\textit{valor}\ \textit{parametros}\mid\varepsilon$ (recursión por la derecha, equivalente) & prods. 23--24 \\
$\textit{condicion}\to\textit{condicion}\ \texttt{OR}\ \textit{termino}\mid\textit{termino}$ &
$\textit{condicion}\to\textit{termino}\ \textit{condicion'}$ & prods. 5--7 \\
 & $\textit{condicion'}\to\texttt{OR}\ \textit{termino}\ \textit{condicion'}\mid\varepsilon$ & \\
$\textit{termino}\to\textit{termino}\ \texttt{AND}\ \textit{factor}\mid\textit{factor}$ &
$\textit{termino}\to\textit{factor}\ \textit{termino'}$ & prods. 8--10 \\
 & $\textit{termino'}\to\texttt{AND}\ \textit{factor}\ \textit{termino'}\mid\varepsilon$ & \\
\bottomrule
\end{tabular}
\end{center}

\paragraph{Verificación en la gramática final.} No hay recursión izquierda directa (lo comprueba una prueba unitaria: ninguna producción
empieza con su propio lado izquierdo). Tampoco indirecta: las dependencias por el primer símbolo forman cadenas sin ciclos:
\begin{itemize}[nosep]
\item $\textit{programa}\to\textit{regla}\to\texttt{WHEN}$;
\item $\textit{condicion}\to\textit{termino}\to\textit{factor}\to\textit{comparacion}\to\texttt{ID}$;
\item $\textit{acciones}\to\textit{accion}\to\texttt{ID}$.
\end{itemize}
La única recursión de \textit{factor} ($\textit{factor}\to\texttt{NOT}\ \textit{factor}$) empieza con un terminal, por lo que es inofensiva.

\subsection{Factorización por la izquierda}
En $G_0$, \textit{programa} tiene dos alternativas con prefijo común \textit{regla}
($\textit{programa}\to\textit{regla}\ \textit{programa}\mid\textit{regla}$): con un solo token de lookahead no se
puede elegir. Se factoriza: $\textit{programa}\to\textit{regla}\ \textit{programa'}$ con
$\textit{programa'}\to\textit{programa}\mid\varepsilon$; como $\textit{programa}$ empieza con $\textit{regla}$, se
sustituye y queda $\textit{programa'}\to\textit{regla}\ \textit{programa'}\mid\varepsilon$ (prods. 1--3).
En el resto de la gramática las alternativas de un mismo no terminal comienzan con terminales distintos
(p.\,ej. \textit{factor}: \texttt{NOT}, \texttt{PARENT\_IZQ}, \texttt{ID}), por lo que \textbf{no se requiere más factorización}; esto se
confirma formalmente en la sección~6 (conjuntos de selección disjuntos).

\subsection{Ambigüedad}
$G_0$ es \textbf{ambigua}: la cadena \texttt{a > 1 OR b > 2 AND c > 3} tiene dos árboles de derivación, según si la
producción $\textit{condicion}\to\textit{condicion}\ \texttt{OR}\ \textit{condicion}$ o
$\textit{condicion}\to\textit{condicion}\ \texttt{AND}\ \textit{condicion}$ se aplica en la raíz:
\begin{center}
\texttt{(a > 1 OR b > 2) AND c > 3} \quad frente a \quad \texttt{a > 1 OR (b > 2 AND c > 3)}.
\end{center}
Además \texttt{a OR b OR c} puede agruparse a la izquierda o a la derecha (ambigüedad de asociatividad). Esto cambiaría
el significado del programa, así que debe eliminarse en la gramática, no en el código.

\paragraph{Solución.} Se separan los \textbf{niveles de precedencia} (\textit{condicion} para OR, \textit{termino} para AND,
\textit{factor} para NOT/paréntesis/comparación): AND liga más fuerte que OR, y NOT más que AND. Con las producciones
5--13 la cadena anterior tiene un único árbol, el de \texttt{a > 1 OR (b > 2 AND c > 3)}. La asociatividad es
\textbf{izquierda}: el AST se construye en un ciclo que agrupa \texttt{((a OR b) OR c)} (sección~7).
Como además la gramática final es LL(1) (sección~6), es \textbf{no ambigua} (toda gramática LL(1) lo es).
La guía del enunciado (\textit{condicion2} con AND/OR mezclados) no define precedencia entre ambos conectores; la estratificación
usada aquí sí lo hace.

Otro punto revisado: en \textit{accion}, los parámetros son identificadores/valores y podrían confundirse con el
siguiente elemento; no hay ambigüedad porque la lista de parámetros termina siempre en \texttt{;} o \texttt{,}
(FOLLOW(parametros) = \{\texttt{PUNTO\_COMA}, \texttt{COMA}\}, ver sección~5).

\subsection{Extensión creativa y su efecto en la gramática}
La extensión de DomoDSL tiene tres componentes, todos integrados en la gramática, FIRST/FOLLOW, tabla, código y pruebas:
\begin{enumerate}[nosep]
\item \textbf{Negación} \texttt{NOT}: $\textit{factor}\to\texttt{NOT}\ \textit{factor}$ (prod.~11); permite \texttt{NOT NOT x == 1}.
\item \textbf{Agrupación con paréntesis}: $\textit{factor}\to\texttt{PARENT\_IZQ}\ \textit{condicion}\ \texttt{PARENT\_DER}$ (prod.~12); es la que introduce
\texttt{PARENT\_DER} en FOLLOW(\textit{condicion}) y justifica los mensajes ``falta \texttt{)}'' / ``\texttt{)} sobrante''.
\item \textbf{Acciones múltiples con parámetros} (propia del dominio): $\textit{acciones}$ con \texttt{COMA} (prods.~19--21) y
\textit{parametros} (prods.~23--24); p.\,ej. \texttt{ajustar termostato 22 TRUE, activar alarma}.
\end{enumerate}

\section{FIRST y FOLLOW}

\paragraph{Definiciones.} $\mathrm{FIRST}(\alpha)$ es el conjunto de terminales que pueden iniciar una cadena derivada de $\alpha$
(incluye $\varepsilon$ si $\alpha\Rightarrow^{*}\varepsilon$). $\mathrm{FOLLOW}(A)$ es el conjunto de terminales que pueden aparecer
inmediatamente a la derecha de $A$ en alguna forma sentencial (incluye \texttt{\$} si $A$ puede ir al final).

\paragraph{Reglas usadas.}
\begin{enumerate}[nosep]
\item Si $X$ es terminal, $\mathrm{FIRST}(X)=\{X\}$. Si $X\to\varepsilon$, $\varepsilon\in\mathrm{FIRST}(X)$.
\item Si $X\to Y_1Y_2\cdots Y_k$: se agrega $\mathrm{FIRST}(Y_1)\setminus\{\varepsilon\}$; si $Y_1$ es anulable, también $\mathrm{FIRST}(Y_2)\setminus\{\varepsilon\}$, y así sucesivamente; $\varepsilon\in\mathrm{FIRST}(X)$ si todos los $Y_i$ son anulables.
\item \texttt{\$}$\in\mathrm{FOLLOW}(S)$.
\item Si $A\to\alpha B\beta$: $\mathrm{FIRST}(\beta)\setminus\{\varepsilon\}\subseteq\mathrm{FOLLOW}(B)$.
\item Si $A\to\alpha B$ o $A\to\alpha B\beta$ con $\varepsilon\in\mathrm{FIRST}(\beta)$: $\mathrm{FOLLOW}(A)\subseteq\mathrm{FOLLOW}(B)$.
\end{enumerate}
Los conjuntos se calculan iterando hasta alcanzar el punto fijo (\texttt{calcular\_first} y \texttt{calcular\_follow} en
\texttt{gramatica.py}); los resultados coinciden con el cálculo manual justificado en las tablas y están verificados por
pruebas unitarias (\texttt{test\_first}, \texttt{test\_follow}).

\subsection{Conjuntos FIRST}
\begin{longtable}{@{}>{\raggedright\arraybackslash}p{2.6cm} >{\raggedright\arraybackslash}p{5.3cm} >{\raggedright\arraybackslash}p{6.9cm}@{}}
\toprule
\textbf{No terminal} & \textbf{FIRST} & \textbf{Justificación} \\
\midrule
\endhead
\textit{programa} & $\{$ \texttt{WHEN} $\}$ & \footnotesize Prod. 1: empieza por \textit{regla}, que no deriva $\varepsilon$; FIRST(programa) = FIRST(regla). \\
\textit{programa}$'$ & $\{$ \texttt{WHEN}, $\varepsilon$ $\}$ & \footnotesize Prod. 2: FIRST(regla) = \{WHEN\}; prod. 3 aporta $\varepsilon$. \\
\textit{regla} & $\{$ \texttt{WHEN} $\}$ & \footnotesize Prod. 4: empieza con el terminal WHEN. \\
\textit{condicion} & $\{$ \texttt{NOT}, \texttt{PARENT\_IZQ}, \texttt{ID} $\}$ & \footnotesize Prod. 5: \textit{termino} no es anulable, luego FIRST(condicion) = FIRST(termino). \\
\textit{condicion}$'$ & $\{$ \texttt{OR}, $\varepsilon$ $\}$ & \footnotesize Prod. 6 empieza con OR; prod. 7 aporta $\varepsilon$. \\
\textit{termino} & $\{$ \texttt{NOT}, \texttt{PARENT\_IZQ}, \texttt{ID} $\}$ & \footnotesize Prod. 8: \textit{factor} no es anulable, luego FIRST(termino) = FIRST(factor). \\
\textit{termino}$'$ & $\{$ \texttt{AND}, $\varepsilon$ $\}$ & \footnotesize Prod. 9 empieza con AND; prod. 10 aporta $\varepsilon$. \\
\textit{factor} & $\{$ \texttt{NOT}, \texttt{PARENT\_IZQ}, \texttt{ID} $\}$ & \footnotesize Prod. 11: NOT; prod. 12: PARENT\_IZQ; prod. 13: FIRST(comparacion) = \{ID\}. \\
\textit{comparacion} & $\{$ \texttt{ID} $\}$ & \footnotesize Prod. 14: empieza con el terminal ID. \\
\textit{valor} & $\{$ \texttt{ID}, \texttt{NUM}, \texttt{TRUE}, \texttt{FALSE} $\}$ & \footnotesize Prods. 15--18: NUM, TRUE, FALSE e ID. \\
\textit{acciones} & $\{$ \texttt{ID} $\}$ & \footnotesize Prod. 19: FIRST(accion) = \{ID\}; \textit{accion} no es anulable. \\
\textit{acciones}$'$ & $\{$ \texttt{COMA}, $\varepsilon$ $\}$ & \footnotesize Prod. 20 empieza con COMA; prod. 21 aporta $\varepsilon$. \\
\textit{accion} & $\{$ \texttt{ID} $\}$ & \footnotesize Prod. 22: empieza con el terminal ID. \\
\textit{parametros} & $\{$ \texttt{ID}, \texttt{NUM}, \texttt{TRUE}, \texttt{FALSE}, $\varepsilon$ $\}$ & \footnotesize Prod. 23: FIRST(valor); prod. 24 aporta $\varepsilon$. \\
\bottomrule
\end{longtable}

Los no terminales anulables son \textit{programa'}, \textit{condicion'}, \textit{termino'}, \textit{acciones'} y \textit{parametros}.

\subsection{Conjuntos FOLLOW}
\begin{longtable}{@{}>{\raggedright\arraybackslash}p{2.6cm} >{\raggedright\arraybackslash}p{5.3cm} >{\raggedright\arraybackslash}p{6.9cm}@{}}
\toprule
\textbf{No terminal} & \textbf{FOLLOW} & \textbf{Justificación} \\
\midrule
\endhead
\textit{programa} & $\{$ \texttt{\$} $\}$ & \footnotesize Símbolo inicial: se agrega \$. \\
\textit{programa}$'$ & $\{$ \texttt{\$} $\}$ & \footnotesize Prod. 1: \textit{programa'} está al final, luego FOLLOW(programa) $\subseteq$ FOLLOW(programa'); la prod. 2 solo repite la regla. \\
\textit{regla} & $\{$ \texttt{WHEN}, \texttt{\$} $\}$ & \footnotesize Prods. 1 y 2: sigue \textit{programa'}; FIRST(programa')$\setminus\{\varepsilon\}$ = \{WHEN\} y como es anulable se agrega FOLLOW(programa') = \{\$\}. \\
\textit{condicion} & $\{$ \texttt{THEN}, \texttt{PARENT\_DER} $\}$ & \footnotesize Prod. 4: sigue THEN. Prod. 12: sigue PARENT\_DER. \\
\textit{condicion}$'$ & $\{$ \texttt{THEN}, \texttt{PARENT\_DER} $\}$ & \footnotesize Prods. 5 y 6: va al final, luego hereda FOLLOW(condicion). \\
\textit{termino} & $\{$ \texttt{THEN}, \texttt{OR}, \texttt{PARENT\_DER} $\}$ & \footnotesize Prods. 5 y 6: sigue \textit{condicion'}; FIRST = \{OR\} y es anulable, luego se agrega FOLLOW(condicion). \\
\textit{termino}$'$ & $\{$ \texttt{THEN}, \texttt{OR}, \texttt{PARENT\_DER} $\}$ & \footnotesize Prods. 8 y 9: va al final, luego hereda FOLLOW(termino). \\
\textit{factor} & $\{$ \texttt{THEN}, \texttt{OR}, \texttt{AND}, \texttt{PARENT\_DER} $\}$ & \footnotesize Prods. 8 y 9: sigue \textit{termino'}: AND, y por ser anulable FOLLOW(termino). Prod. 11 (NOT factor): hereda FOLLOW(factor) (sin efecto nuevo). \\
\textit{comparacion} & $\{$ \texttt{THEN}, \texttt{OR}, \texttt{AND}, \texttt{PARENT\_DER} $\}$ & \footnotesize Prod. 13: va al final, luego hereda FOLLOW(factor). \\
\textit{valor} & $\{$ \texttt{THEN}, \texttt{PUNTO\_COMA}, \texttt{OR}, \texttt{AND}, \texttt{PARENT\_DER}, \texttt{ID}, \texttt{NUM}, \texttt{TRUE}, \texttt{FALSE}, \texttt{COMA} $\}$ & \footnotesize Prod. 14: va al final, hereda FOLLOW(comparacion). Prod. 23: sigue \textit{parametros}: FIRST(parametros)$\setminus\{\varepsilon\}$ y, al ser anulable, FOLLOW(parametros). \\
\textit{acciones} & $\{$ \texttt{PUNTO\_COMA} $\}$ & \footnotesize Prod. 4: sigue PUNTO\_COMA. \\
\textit{acciones}$'$ & $\{$ \texttt{PUNTO\_COMA} $\}$ & \footnotesize Prods. 19 y 20: va al final, luego hereda FOLLOW(acciones). \\
\textit{accion} & $\{$ \texttt{PUNTO\_COMA}, \texttt{COMA} $\}$ & \footnotesize Prods. 19 y 20: sigue \textit{acciones'}: COMA, y por ser anulable FOLLOW(acciones) (y FOLLOW(acciones')). \\
\textit{parametros} & $\{$ \texttt{PUNTO\_COMA}, \texttt{COMA} $\}$ & \footnotesize Prod. 22: va al final, hereda FOLLOW(accion). Prod. 23: va al final, hereda FOLLOW(parametros). \\
\bottomrule
\end{longtable}

\section{Tabla LL(1)}

\paragraph{Construcción.} Para cada producción $A\to\alpha$ (numerada $n$): (a) para cada terminal $a\in\mathrm{FIRST}(\alpha)\setminus\{\varepsilon\}$,
se escribe $n$ en $M[A,a]$; (b) si $\varepsilon\in\mathrm{FIRST}(\alpha)$, se escribe $n$ en $M[A,b]$ para cada
$b\in\mathrm{FOLLOW}(A)$ (incluido \texttt{\$}). Una celda vacía significa error sintáctico.

\begin{landscape}
\begin{table}[H]
\centering
\caption{Tabla predictiva LL(1) de DomoDSL (cada número es la producción de la sección 3; celda vacía = error).}
\normalsize\setlength{\tabcolsep}{2pt}\renewcommand{\arraystretch}{1.25}
\begin{tabular}{@{}l|>{\centering\arraybackslash}p{1.4cm}>{\centering\arraybackslash}p{1.4cm}>{\centering\arraybackslash}p{1.4cm}>{\centering\arraybackslash}p{1.4cm}>{\centering\arraybackslash}p{1.4cm}>{\centering\arraybackslash}p{1.4cm}>{\centering\arraybackslash}p{1.4cm}>{\centering\arraybackslash}p{1.4cm}>{\centering\arraybackslash}p{1.4cm}>{\centering\arraybackslash}p{1.4cm}>{\centering\arraybackslash}p{1.4cm}>{\centering\arraybackslash}p{1.4cm}>{\centering\arraybackslash}p{1.4cm}>{\centering\arraybackslash}p{1.4cm}>{\centering\arraybackslash}p{1.4cm}@{}}
\toprule
 & {\fontsize{6.5}{7}\selectfont \texttt{WHEN}} & {\fontsize{6.5}{7}\selectfont \texttt{THEN}} & {\fontsize{6.5}{7}\selectfont \texttt{PUNTO\_COMA}} & {\fontsize{6.5}{7}\selectfont \texttt{OR}} & {\fontsize{6.5}{7}\selectfont \texttt{AND}} & {\fontsize{6.5}{7}\selectfont \texttt{NOT}} & {\fontsize{6.5}{7}\selectfont \texttt{PARENT\_IZQ}} & {\fontsize{6.5}{7}\selectfont \texttt{PARENT\_DER}} & {\fontsize{6.5}{7}\selectfont \texttt{ID}} & {\fontsize{6.5}{7}\selectfont \texttt{OP\_REL}} & {\fontsize{6.5}{7}\selectfont \texttt{NUM}} & {\fontsize{6.5}{7}\selectfont \texttt{TRUE}} & {\fontsize{6.5}{7}\selectfont \texttt{FALSE}} & {\fontsize{6.5}{7}\selectfont \texttt{COMA}} & {\fontsize{6.5}{7}\selectfont \texttt{\$}} \\
\midrule
\textit{programa} & \cellcolor{celda}\textbf{1} &  &  &  &  &  &  &  &  &  &  &  &  &  &  \\
\textit{programa}$'$ & \cellcolor{celda}\textbf{2} &  &  &  &  &  &  &  &  &  &  &  &  &  & \cellcolor{celda}\textbf{3} \\
\textit{regla} & \cellcolor{celda}\textbf{4} &  &  &  &  &  &  &  &  &  &  &  &  &  &  \\
\textit{condicion} &  &  &  &  &  & \cellcolor{celda}\textbf{5} & \cellcolor{celda}\textbf{5} &  & \cellcolor{celda}\textbf{5} &  &  &  &  &  &  \\
\textit{condicion}$'$ &  & \cellcolor{celda}\textbf{7} &  & \cellcolor{celda}\textbf{6} &  &  &  & \cellcolor{celda}\textbf{7} &  &  &  &  &  &  &  \\
\textit{termino} &  &  &  &  &  & \cellcolor{celda}\textbf{8} & \cellcolor{celda}\textbf{8} &  & \cellcolor{celda}\textbf{8} &  &  &  &  &  &  \\
\textit{termino}$'$ &  & \cellcolor{celda}\textbf{10} &  & \cellcolor{celda}\textbf{10} & \cellcolor{celda}\textbf{9} &  &  & \cellcolor{celda}\textbf{10} &  &  &  &  &  &  &  \\
\textit{factor} &  &  &  &  &  & \cellcolor{celda}\textbf{11} & \cellcolor{celda}\textbf{12} &  & \cellcolor{celda}\textbf{13} &  &  &  &  &  &  \\
\textit{comparacion} &  &  &  &  &  &  &  &  & \cellcolor{celda}\textbf{14} &  &  &  &  &  &  \\
\textit{valor} &  &  &  &  &  &  &  &  & \cellcolor{celda}\textbf{18} &  & \cellcolor{celda}\textbf{15} & \cellcolor{celda}\textbf{16} & \cellcolor{celda}\textbf{17} &  &  \\
\textit{acciones} &  &  &  &  &  &  &  &  & \cellcolor{celda}\textbf{19} &  &  &  &  &  &  \\
\textit{acciones}$'$ &  &  & \cellcolor{celda}\textbf{21} &  &  &  &  &  &  &  &  &  &  & \cellcolor{celda}\textbf{20} &  \\
\textit{accion} &  &  &  &  &  &  &  &  & \cellcolor{celda}\textbf{22} &  &  &  &  &  &  \\
\textit{parametros} &  &  & \cellcolor{celda}\textbf{24} &  &  &  &  &  & \cellcolor{celda}\textbf{23} &  & \cellcolor{celda}\textbf{23} & \cellcolor{celda}\textbf{23} & \cellcolor{celda}\textbf{23} & \cellcolor{celda}\textbf{24} &  \\
\bottomrule
\end{tabular}
\end{table}
\end{landscape}

\subsection{¿Es la gramática LL(1)?}
\textbf{Sí.} Ninguna celda contiene más de una producción (el programa lo verifica: \texttt{CONFLICTOS = []}).
Equivalentemente, para cada no terminal con varias alternativas los \emph{conjuntos de selección}
$\mathrm{SEL}(A\to\alpha)=\mathrm{FIRST}(\alpha)\setminus\{\varepsilon\}\ \cup\ (\mathrm{FOLLOW}(A)\text{ si }\varepsilon\in\mathrm{FIRST}(\alpha))$
son disjuntos de a pares:

\begin{longtable}{@{}>{\raggedright\arraybackslash}p{2.3cm} >{\raggedright\arraybackslash}p{4.0cm} >{\raggedright\arraybackslash}p{5.9cm} c@{}}
\toprule
\textbf{No terminal} & \textbf{Producción} & \textbf{Conjunto de selección} & \textbf{¿Disjuntos?} \\
\midrule
\endhead
\textit{programa}$'$ & (2)  \textit{regla} \textit{programa}$'$ & $\{$ \texttt{WHEN} $\}$ & \textbf{sí} \\
 & (3)  $\varepsilon$ & $\{$ \texttt{\$} $\}$ &  \\
\midrule
\textit{condicion}$'$ & (6)  \texttt{OR} \textit{termino} \textit{condicion}$'$ & $\{$ \texttt{OR} $\}$ & \textbf{sí} \\
 & (7)  $\varepsilon$ & $\{$ \texttt{THEN}, \texttt{PARENT\_DER} $\}$ &  \\
\midrule
\textit{termino}$'$ & (9)  \texttt{AND} \textit{factor} \textit{termino}$'$ & $\{$ \texttt{AND} $\}$ & \textbf{sí} \\
 & (10)  $\varepsilon$ & $\{$ \texttt{THEN}, \texttt{OR}, \texttt{PARENT\_DER} $\}$ &  \\
\midrule
\textit{factor} & (11)  \texttt{NOT} \textit{factor} & $\{$ \texttt{NOT} $\}$ & \textbf{sí} \\
 & (12)  \texttt{PARENT\_IZQ} \textit{condicion} \texttt{PARENT\_DER} & $\{$ \texttt{PARENT\_IZQ} $\}$ &  \\
 & (13)  \textit{comparacion} & $\{$ \texttt{ID} $\}$ &  \\
\midrule
\textit{valor} & (15)  \texttt{NUM} & $\{$ \texttt{NUM} $\}$ & \textbf{sí} \\
 & (16)  \texttt{TRUE} & $\{$ \texttt{TRUE} $\}$ &  \\
 & (17)  \texttt{FALSE} & $\{$ \texttt{FALSE} $\}$ &  \\
 & (18)  \texttt{ID} & $\{$ \texttt{ID} $\}$ &  \\
\midrule
\textit{acciones}$'$ & (20)  \texttt{COMA} \textit{accion} \textit{acciones}$'$ & $\{$ \texttt{COMA} $\}$ & \textbf{sí} \\
 & (21)  $\varepsilon$ & $\{$ \texttt{PUNTO\_COMA} $\}$ &  \\
\midrule
\textit{parametros} & (23)  \textit{valor} \textit{parametros} & $\{$ \texttt{ID}, \texttt{NUM}, \texttt{TRUE}, \texttt{FALSE} $\}$ & \textbf{sí} \\
 & (24)  $\varepsilon$ & $\{$ \texttt{PUNTO\_COMA}, \texttt{COMA} $\}$ &  \\
\bottomrule
\end{longtable}

\paragraph{Puntos delicados y por qué no hay conflicto.}
\begin{itemize}[nosep]
\item \textit{parametros}: las producciones $\textit{valor}\ \textit{parametros}$ (selección $\{\texttt{NUM},\texttt{TRUE},\texttt{FALSE},\texttt{ID}\}$) y $\varepsilon$ (selección
$\mathrm{FOLLOW}=\{\texttt{PUNTO\_COMA},\texttt{COMA}\}$) no se solapan, precisamente porque la lista de parámetros siempre termina en
\texttt{;} o \texttt{,}.
\item \textit{termino'} y \textit{condicion'}: la alternativa vacía se elige con $\mathrm{FOLLOW}$ (\texttt{OR}, \texttt{THEN}, \texttt{)}), que no
contiene \texttt{AND}; por eso AND se asocia siempre antes que OR (precedencia).
\item \textit{programa'}: \texttt{WHEN} inicia otra regla y \texttt{\$} termina el programa.
\end{itemize}
\paragraph{Conclusión.} No fue necesario modificar la gramática después de calcular la tabla: las transformaciones de la
sección~4 ya dejaron una gramática LL(1) utilizable por ambos analizadores.

